\documentclass[11pt]{article}
\usepackage{ochamath}
\subjectcolor{2F5DA8}
\setsheet{線形代数}{逆行列と正則性　演習}{https://university-math-crj.pages.dev/linear-algebra/inverse-matrices/}
\begin{document}
\makesheethead{解答}

\problem[逆行列と復元]
\begin{ans}
$3\times1-1\times2=1$ より \[A^{-1}=\begin{pmatrix}1&-1\\-2&3\end{pmatrix},\qquad \boldsymbol x=\begin{pmatrix}7-5\\-14+15\end{pmatrix}=\begin{pmatrix}2\\1\end{pmatrix}.\] $A(2,1)^{\mathsf T}=(7,5)^{\mathsf T}$ で検算。
\end{ans}
\point{逆行列そのものも積がIになることで確かめる。}

\problem[正則でない変換]
\begin{ans}
$\boldsymbol v=(-2,1)^{\mathsf T}$ は非零だが $B\boldsymbol v=(0,0)^{\mathsf T}$。もし逆行列があれば $\boldsymbol v=B^{-1}\boldsymbol0=\boldsymbol0$ となり矛盾する。
\end{ans}
\point{情報を失う方向があると逆変換は存在しない。}
\newpage

\problem[片側逆]
\begin{ans}
$LA=I_2$、$AL=\operatorname{diag}(1,1,0)$ で $I_3$ ではない。$L$ は左逆行列だが通常の逆行列ではない。第3成分を0にする変換が残る。
\end{ans}
\point{長方形行列では左右の積のサイズも異なる。}

\problem[センサ校正]
\begin{ans}
$C^{-1}=C/2$ より $x_1=(6+2)/2=4,\ x_2=(6-2)/2=2$。検算：$4+2=6,\ 4-2=2$。
\end{ans}
\point{差の出力があるため2つの入力を分離できる。}
\end{document}
